\clearpage
\SourceLeaf{70}
\pdfbookmark[1]{Observationes cometae secundi a. 1813. Anzeige von 1814}{anzeige-cometae1814}
\begin{center}\rule{.75\textwidth}{.3pt}\\
Göttingische gelehrte Anzeigen. 3. Januar 1814.\\
\rule{.75\textwidth}{.3pt}\end{center}
Wir haben noch die Anzeige einer kurzen Abhandlung nachzuholen, die Professor \textsc{Gauß} am 10. September des vergangenen Jahres der Königlichen Gesellschaft vorlegte. Sie trägt den Titel
\begin{center}
\textit{Observationes cometae secundi a.} 1813 \textit{in observatorio Gottingensi factae, adiectis
nonnullis adnotationibus circa calculum orbitarum parabolicarum.}
\end{center}
Bekanntlich entdeckte unser Professor \textsc{Harding} diesen Kometen am 3. April des vergangenen Jahres im Poniatowskischen Stier. Die an der hiesigen Sternwarte angestellten Beobachtungen reichen vom 9. bis zum 25. April. In der Abhandlung findet man sie \textit{vollständig}. Da der größte Teil schon andernorts bekannt gemacht worden ist, übergehen wir sie hier mit Stillschweigen und geben aus der Abhandlung nur die parabolische Bahn an, die Professor \textsc{Gauß} von Dr. \textsc{Gerling} in Cassel aus den hiesigen und einigen Bremer und Pariser Beobachtungen berechnen ließ.
\begin{center}\small
Zeit des Durchgangs durch das Perihel, 19. Mai 1813, $10^{\mathrm h}\;24^{\mathrm m}\;5^{\mathrm s}$\\
Mittlere Göttinger Zeit
\end{center}
\begin{center}\small
\begin{tabular}{@{}l@{\qquad}r@{}}
Länge des Perihels $\quad.\;.\;.\;.$ & $197^0\;43'\;7''7$\\
Länge des aufsteigenden Knotens $\quad,$ & $42\;40\;15.2$\\
Bahnneigung $\quad.\;.\;.\;.\;.\;.$ & $81\quad2\;11.8$\\
Logarithmus des kleinsten Abstands & $0.0849212$\\
\multicolumn{2}{l}{Bewegung \textit{rückläufig.}}
\end{tabular}\end{center}

\SourceLeaf{71}\SourceGathering{VI.~9}%
Noch vor wenigen Jahrzehnten konnten in ganz Europa nur wenige Menschen eine Kometenbahn berechnen. Heute haben verbesserte Methoden diese Arbeit so sehr erleichtert und vereinfacht, dass sich ein sonst fähiger Kopf ohne Schwierigkeiten mit ihr vertraut machen und eine Kometenbahn in weniger Stunden bestimmen kann, als früher Tage nötig waren. Für die \textit{erste} Berechnung einer Bahn aus drei Beobachtungen lässt die bequeme Methode von Dr. \textsc{Olbers} fast nichts zu wünschen übrig, sofern man Beobachtungen wählen kann, bei denen die Richtung der geozentrischen Bewegung der Richtung des Großkreises nicht zu nahe kommt, der auf der Himmelskugel durch den mittleren geozentrischen Ort des Kometen und den entsprechenden Sonnenort gezogen wird. Der zweite Teil der vorliegenden Abhandlung soll einige Punkte dieser schönen Methode noch etwas vereinfachen und für die numerische Rechnung noch bequemer machen. Der Versuch betrifft das indirekte Verfahren zur näherungsweisen Bestimmung der Abstände des Kometen von der Sonne und der Sehne zwischen den beiden äußersten Orten. Nur wer in der astronomischen Rechnung unerfahren ist, könnte es der Methode zum Vorwurf machen, dass dieses Verfahren indirekt ist. Die Form, die ihm \textsc{Olbers} gegeben hat, hat jedoch einen Nachteil. In den meisten Fällen würden die kleineren Tafeln mit fünf Dezimalstellen zu große Fehler verursachen, auch wenn man, wie bei einer ersten Näherung angemessen, in der Rechnung keine äußerste Genauigkeit anstrebt. Das liegt daran, dass mehrere ihrer Natur nach sehr kleine Größen als Differenzen anderer, fast gleicher Größen erscheinen. Professor \textsc{Gauß} suchte diesem Nachteil zu begegnen, indem er mehrere Hilfsgrößen einführte, die sich mit den kleinen Tafeln immer hinreichend genau bestimmen lassen. Dadurch gab er den Formeln für die Abstände von der Sonne und für die Sehne eine andere Form, bei der auch die anschließenden Versuche der indirekten Lösung etwas bequemer zu werden scheinen. Diese Umformungen lassen sich hier nicht im Auszug wiedergeben, ohne dass wir diesen Teil der Abhandlung ganz abschreiben. Wir verweisen daher auf den demnächst erscheinenden zweiten Band der \textit{Commentationes} der Gesellschaft, in dem die ganze Abhandlung bereits abgedruckt ist.

Der Verfasser glaubte manchen Lesern einen Dienst zu erweisen, wenn er die ihm eigenen Änderungen mit
